\documentclass[11pt,a4paper]{article}

% --- Page layout ---
\usepackage[margin=1in]{geometry}
\usepackage{parskip}

% --- Fonts and encoding ---
\usepackage[T1]{fontenc}
\usepackage[utf8]{inputenc}
\usepackage{lmodern}
\usepackage{microtype}

% --- Math ---
\usepackage{amsmath,amssymb}

% --- Graphics and color ---
\usepackage{graphicx}
\usepackage[dvipsnames,table]{xcolor}

% --- Tables ---
\usepackage{booktabs}
\usepackage{tabularx}
\usepackage{multirow}

% --- Code listings ---
\usepackage{listings}
\lstset{
  basicstyle=\ttfamily\small,
  keywordstyle=\color{MidnightBlue}\bfseries,
  commentstyle=\color{OliveGreen},
  stringstyle=\color{BrickRed},
  breaklines=true,
  frame=single,
  rulecolor=\color{gray!40},
  numbers=left,
  numberstyle=\tiny\color{gray},
  xleftmargin=2em,
  framexleftmargin=1.5em,
  columns=flexible,
  showstringspaces=false,
}

% --- Hyperlinks ---
\usepackage[colorlinks=true,linkcolor=MidnightBlue,citecolor=OliveGreen,urlcolor=BrickRed]{hyperref}

% --- Bibliography ---
\usepackage[numbers,sort&compress]{natbib}

% --- Section styling ---
\usepackage{titlesec}
\titleformat{\section}{\large\bfseries\color{MidnightBlue}}{\thesection}{1em}{}
\titleformat{\subsection}{\normalsize\bfseries\color{MidnightBlue!80}}{\thesubsection}{1em}{}

% --- Colored table header helper ---
\newcommand{\thead}[1]{\textbf{\color{white}#1}}
\definecolor{HeaderBlue}{HTML}{2C3E6B}

% --- Document ---
\begin{document}

% --- Title block ---
\begin{center}
  {\LARGE\bfseries rvLLM: A High-Performance LLM Inference Engine in Rust\par}
  \vspace{0.6em}
  {\large m0at\par}
  \vspace{0.3em}
  {\normalsize\texttt{https://github.com/m0at/rvllm}\par}
  \vspace{0.3em}
  {\normalsize March 2026\par}
\end{center}
\vspace{1em}

% --- Abstract ---
\begin{abstract}
We present rvLLM, a from-scratch Rust implementation of the vLLM inference engine
targeting OpenAI-compatible API serving with CUDA GPU acceleration. The system
comprises 22~Rust crates and pre-compiled CUDA PTX kernels, compiling to a single
15\,MB static binary. We describe the modular crate architecture, CUDA integration
strategy via PTX kernels loaded through \texttt{cudarc}, and key design decisions
including a mock-GPU feature for hardware-free development, CPU attention fallback,
and PyTorch weight-layout compatibility. The merged main branch produces coherent
inference output on NVIDIA A100 with a $\sim$7\,s startup time and 333\,MB resident
memory. Full throughput benchmarks against Python vLLM are being refreshed in a
follow-up pass. rvLLM is open-source under the Apache-2.0 license.
\end{abstract}

% ============================================================================
\section{Introduction}
\label{sec:introduction}
% ============================================================================

Large language model (LLM) inference serving has become a critical infrastructure
component. The dominant open-source serving engine, vLLM~\cite{kwon2023vllm},
introduced PagedAttention for efficient KV cache management and demonstrated that
virtual-memory-inspired techniques can dramatically improve GPU memory utilisation
during batched inference. However, vLLM is implemented in Python with PyTorch as its
tensor runtime, which introduces several practical limitations.

First, \textbf{Python's Global Interpreter Lock} (GIL) serialises the scheduler,
tokenizer, and output-processing stages. At high concurrency (hundreds of concurrent
sequences), the scheduling loop itself becomes a throughput bottleneck---not the GPU
computation.

Second, \textbf{garbage collection pauses} grow as Python tracks millions of tensor
metadata objects during large-batch serving. These pauses introduce tail-latency
spikes that are difficult to diagnose or eliminate.

Third, \textbf{deployment complexity} is substantial. A fresh \texttt{pip install vllm}
pulls approximately 500\,MB of dependencies including PyTorch ($\sim$2\,GB with CUDA
runtime), the \texttt{transformers} library, NumPy, and dozens of supporting packages.
Startup requires 30--60 seconds for Python imports, Triton JIT compilation, and NCCL
initialisation.

rvLLM addresses these limitations by reimplementing the core vLLM architecture in Rust.
The system compiles to a single 14\,MB statically-linked binary with zero runtime
dependencies beyond the CUDA driver. Rust's ownership model provides deterministic
memory deallocation without garbage collection, and the absence of a GIL enables true
parallelism across the scheduler, tokenizer, and sampling stages. Direct cuBLAS and
CUDA kernel invocation through \texttt{cudarc}~\cite{cudarc2023} eliminates the
PyTorch middleware layer.

The primary contributions of this work are:
\begin{enumerate}
  \item A modular 22-crate Rust workspace architecture for LLM inference serving,
        with clear separation between GPU primitives, model execution, scheduling,
        and API layers.
  \item A CUDA integration strategy using 13~pre-compiled PTX kernels loaded at
        runtime, avoiding JIT compilation and enabling deterministic startup.
  \item A mock-GPU feature that allows the entire scheduling, sampling, and API
        stack (769~tests across 22~crates) to be developed and tested without GPU
        hardware.
  \item Preliminary benchmark results demonstrating competitive throughput on A100
        GPUs with substantially reduced memory footprint, binary size, and startup
        time.
\end{enumerate}

% ============================================================================
\section{Architecture}
\label{sec:architecture}
% ============================================================================

rvLLM is organised as a Cargo workspace containing 22~member crates plus a server
binary crate. Each crate has a single well-defined responsibility, and dependencies
flow strictly downward from the API layer to the GPU primitives. We describe the key
layers below.

\subsection{Core and Configuration}

\texttt{rvllm-core} provides shared types, a structured error hierarchy using
\texttt{thiserror}, and a prelude module. \texttt{rvllm-config} handles CLI argument
parsing (via \texttt{clap}), TOML configuration files, and validation of runtime
parameters such as GPU memory utilisation fraction, maximum sequence length, and
tensor-parallel size. \texttt{rvllm-sequence} defines the sequence state machine
(pending, running, finished, swapped) and request-group metadata.

\subsection{GPU Abstraction Layer}

\texttt{rvllm-gpu} provides a uniform interface over CUDA and a mock backend,
controlled by Cargo feature flags. The CUDA path wraps \texttt{cudarc} for device
management, stream synchronisation, and memory allocation. A \texttt{cuBLAS} wrapper
maps row-major Rust tensors to cuBLAS's column-major convention using the transpose
trick $C = (B^T A^T)^T$, exposing a row-major GEMM interface. The PTX kernel loader
discovers and loads pre-compiled kernels at startup. The mock backend substitutes
CPU-based operations with matching signatures, enabling the full test suite to run on
machines without NVIDIA GPUs.

\subsection{Model Execution}

\texttt{rvllm-model-loader} handles SafeTensors~\cite{safetensors2023} weight loading
from HuggingFace Hub or local paths, including sharded checkpoints and GGUF format.
Weight tensors are memory-mapped for efficient loading.

\texttt{rvllm-model-runner} implements the transformer forward pass. Ten model
architectures are supported (Table~\ref{tab:architectures}), each implementing an
\texttt{Architecture} trait that defines layer structure, attention configuration
(multi-head, grouped-query~\cite{ainslie2023gqa}, or multi-query), and normalization
type (RMSNorm~\cite{zhang2019root} or LayerNorm).

\texttt{rvllm-attention} provides FlashAttention~\cite{dao2022flashattention,
dao2023flashattention2} and PagedAttention~\cite{kwon2023vllm} backends. The CPU
reference implementation enables correctness testing; the CUDA kernels handle
production workloads.

\texttt{rvllm-kv-cache} manages the paged key-value cache with fixed-size blocks,
block tables for virtual-to-physical mapping, and copy-on-write semantics for beam
search.

\texttt{rvllm-quant} supports GPTQ, AWQ, and FP8~\cite{holkman2023fp8} (E4M3)
dequantization with per-head scaling for the KV cache.

\begin{table}[t]
  \centering
  \caption{Supported model architectures in rvLLM~v0.1.0.}
  \label{tab:architectures}
  \small
  \begin{tabular}{@{}ll@{}}
    \toprule
    \rowcolor{HeaderBlue} \thead{Architecture} & \thead{Models} \\
    \midrule
    LlamaForCausalLM      & Llama~2/3, CodeLlama, Vicuna \\
    \rowcolor{gray!8}
    MistralForCausalLM     & Mistral~7B, Mistral Nemo \\
    Qwen2ForCausalLM       & Qwen2, Qwen2.5 \\
    \rowcolor{gray!8}
    PhiForCausalLM         & Phi-2, Phi-3, Phi-3.5 \\
    GemmaForCausalLM       & Gemma, Gemma~2 \\
    \rowcolor{gray!8}
    MixtralForCausalLM     & Mixtral 8x7B, 8x22B (MoE) \\
    DeepseekV2ForCausalLM  & DeepSeek-V2, DeepSeek-V2.5 (MoE) \\
    \rowcolor{gray!8}
    GPTNeoXForCausalLM     & Pythia, GPT-NeoX-20B \\
    StableLmForCausalLM    & StableLM-3B, StableLM-2 \\
    \rowcolor{gray!8}
    CohereForCausalLM      & Command-R, Command-R+ \\
    \bottomrule
  \end{tabular}
\end{table}

\subsection{Scheduling and Execution}

\texttt{rvllm-scheduler} implements continuous batching with first-come-first-served,
priority, and shortest-job-first scheduling policies. Preemption allows long-running
sequences to be swapped to CPU memory when GPU blocks are exhausted, and prefix caching
with LRU eviction avoids redundant KV cache computation for shared prompt prefixes.

\texttt{rvllm-executor} orchestrates single- or multi-GPU execution, dispatching work
to \texttt{rvllm-worker} instances. Each worker owns a GPU device, runs the forward
pass, and returns logits for sampling. \texttt{rvllm-speculative} implements draft-model
speculative decoding~\cite{leviathan2023speculative} as an optional execution strategy.

\texttt{rvllm-engine} ties scheduling and execution together in a Tokio-based async
inference loop. Incoming requests are tokenised, scheduled into batches, executed on
GPU, sampled, and detokenised, with results streamed back via Server-Sent Events (SSE).

\subsection{Sampling}

\texttt{rvllm-sampling} implements the full sampling pipeline: temperature scaling,
top-$k$ filtering, nucleus (top-$p$) sampling, min-$p$ filtering, repetition penalty,
frequency penalty, presence penalty, and multinomial sampling with optional seed-based
determinism. Batch sampling across sequences uses Rayon~\cite{rayon2023} for CPU
parallelism, distributing work across all available cores.

\subsection{API and Observability}

\texttt{rvllm-api} provides an OpenAI-compatible~\cite{openai2023api} HTTP server
built on \texttt{axum}~\cite{axum2023}. Supported endpoints include
\texttt{/v1/completions}, \texttt{/v1/chat/completions}, \texttt{/v1/models},
\texttt{/v1/embeddings}, and \texttt{/v1/batches}, with both streaming (SSE) and
non-streaming response modes. Guided decoding supports JSON mode, JSON schema
validation, and regex grammars. Tool/function calling follows the Hermes convention.

\texttt{rvllm-telemetry} exports Prometheus metrics (forward time, time-to-first-token,
inter-token latency, queue depth) and structured tracing via the \texttt{tracing}
ecosystem.

\texttt{rvllm-tokenizer} wraps HuggingFace's \texttt{tokenizers}
library for encoding, decoding, and chat template rendering.

\texttt{rvllm-python} provides PyO3 bindings, exposing the sampling pipeline and
tokenizer to Python users via \texttt{import rvllm}.

% ============================================================================
\section{CUDA Integration}
\label{sec:cuda}
% ============================================================================

rvLLM interfaces with NVIDIA GPUs through two mechanisms: cuBLAS for dense
matrix operations (GEMM) and custom CUDA kernels for element-wise and reduction
operations. All GPU interaction is mediated through \texttt{cudarc}, which provides
safe Rust bindings to the CUDA driver API.

\subsection{Pre-compiled PTX Kernels}

The system includes 13~hand-written CUDA kernels (Table~\ref{tab:kernels}),
totalling approximately 1,500 lines of CUDA code. Kernels are compiled to PTX
ahead-of-time via \texttt{nvcc}, supporting compute capabilities from \texttt{sm\_70}
(V100) through \texttt{sm\_100} (B100/B200). At server startup, the PTX kernel loader
discovers available \texttt{.ptx} files and registers them with the CUDA driver. This
avoids runtime JIT compilation entirely, contributing to rvLLM's fast startup time.

\begin{table}[t]
  \centering
  \caption{Custom CUDA kernels in rvLLM. All kernels are compiled to PTX for compute capabilities sm\_70 through sm\_100.}
  \label{tab:kernels}
  \small
  \begin{tabular}{@{}llp{5.5cm}@{}}
    \toprule
    \rowcolor{HeaderBlue} \thead{Kernel} & \thead{File} & \thead{Purpose} \\
    \midrule
    PagedAttention V2 & \texttt{paged\_attention.cu} & Block-table indirection with online softmax \\
    \rowcolor{gray!8}
    FlashAttention    & \texttt{flash\_attention.cu} & Tiled attention with shared-memory reduction \\
    RMSNorm (FP32)    & \texttt{rms\_norm.cu}        & Shared-memory parallel reduction \\
    \rowcolor{gray!8}
    RMSNorm (FP16)    & \texttt{rms\_norm\_f16.cu}    & Half-precision variant \\
    Fused Residual+RMSNorm & \texttt{fused\_residual\_rmsnorm.cu} & Fused residual add and normalisation \\
    \rowcolor{gray!8}
    Rotary Embedding  & \texttt{rotary\_embedding.cu} & RoPE with GQA support \\
    Activation (FP32) & \texttt{activation.cu}       & SiLU, GELU, fused SiLU$\times$mul for MLP \\
    \rowcolor{gray!8}
    Activation (FP16) & \texttt{activation\_f16.cu}   & Half-precision variants \\
    Softmax           & \texttt{softmax.cu}          & Warp-level numerically stable softmax \\
    \rowcolor{gray!8}
    Copy Blocks       & \texttt{copy\_blocks.cu}      & KV cache block copy for beam search \\
    Embedding Gather  & \texttt{embedding\_gather.cu}  & Token embedding table lookup \\
    \rowcolor{gray!8}
    Add Bias          & \texttt{add\_bias.cu}         & Fused bias addition \\
    FP8 KV            & \texttt{fp8\_kv.cu}           & E4M3 quantisation with per-head scaling \\
    \bottomrule
  \end{tabular}
\end{table}

\subsection{cuBLAS GEMM Wrapper}

Dense matrix multiplications (linear projections, MLP layers) use cuBLAS SGEMM.
Because cuBLAS assumes column-major layout while Rust arrays are row-major, rvLLM
uses the identity $(AB)^T = B^T A^T$ to compute row-major results. The wrapper
function accepts row-major matrices $A \in \mathbb{R}^{m \times k}$ and
$B \in \mathbb{R}^{k \times n}$, passes them transposed to cuBLAS, and the
result $C \in \mathbb{R}^{m \times n}$ is returned in row-major layout without
explicit transposition. This maintains compatibility with PyTorch's weight layout
convention where linear layer weights are stored as $[\text{out\_features},
\text{in\_features}]$.

\subsection{Memory Management}

GPU memory is managed through explicit allocation pools in \texttt{rvllm-memory}.
The block manager (\texttt{rvllm-block-manager}) maintains a free list of fixed-size
KV cache blocks and supports copy-on-write for beam search and prefix sharing. A swap
manager handles GPU$\leftrightarrow$CPU memory transfers when the block allocator is
exhausted. Pinned (page-locked) host memory is used for efficient DMA transfers.

% ============================================================================
\section{Key Design Decisions}
\label{sec:design}
% ============================================================================

\subsection{Mock-GPU Feature}

A central design decision is the \texttt{mock-gpu} Cargo feature, which replaces all
CUDA calls with CPU-based stubs that produce matching output shapes. This enables:
\begin{itemize}
  \item Development on macOS, Windows, or any Linux machine without an NVIDIA GPU.
  \item The full test suite (769~tests across 22~crates) to run in CI without GPU
        instances.
  \item Independent testing of scheduling, sampling, tokenization, and API logic
        against known outputs.
\end{itemize}
Only the GPU forward pass and CUDA kernel behaviour require real hardware to validate.
All other layers---including the continuous batching scheduler, all sampling algorithms,
the full HTTP API, and streaming---are tested under mock-GPU.

\subsection{CPU Attention Fallback}

rvLLM includes a CPU reference implementation of both standard multi-head attention and
FlashAttention-2. This serves dual purposes: correctness testing (comparing CPU and GPU
outputs token-by-token) and enabling a functional inference path on CPU-only machines.
GPU projections (Q/K/V linear layers via cuBLAS) can be combined with CPU attention
computation in a hybrid configuration.

\subsection{Feature-Gated CUDA Graphs}

CUDA graph capture and replay is implemented but disabled by default behind a Cargo
feature flag. CUDA graphs eliminate kernel launch overhead by recording a sequence of
GPU operations and replaying them in a single driver call. However, graph capture
imposes constraints on dynamic shapes and memory allocation patterns that interact
poorly with the variable batch sizes inherent to continuous batching. rvLLM provides
the infrastructure for CUDA graph capture/replay but defaults to eager execution,
which has proven sufficient at current scale. This is an area of active development.

\subsection{PyTorch Weight Layout Compatibility}

rvLLM loads weights directly from SafeTensors checkpoints produced by PyTorch-based
training pipelines. PyTorch stores linear layer weights in $[\text{out\_features},
\text{in\_features}]$ layout. Rather than transposing weights at load time (which
would require additional memory and time), the cuBLAS wrapper's row-major-to-column-major
mapping naturally accommodates this convention. This means any HuggingFace model
checkpoint can be loaded without weight preprocessing.

\subsection{Deterministic Startup}

rvLLM deliberately avoids runtime code generation. Triton kernels (used by Python
vLLM) require a Python-based JIT compiler that adds 10--30 seconds to startup. By
pre-compiling all kernels to PTX at build time, rvLLM achieves a startup sequence of:
(1) load binary, (2) parse configuration, (3) load PTX kernels, (4) memory-map model
weights, (5) begin serving. Measured startup time is approximately 8~seconds on A100,
dominated by weight loading from disk.

% ============================================================================
\section{Preliminary Results}
\label{sec:results}
% ============================================================================

\textbf{Disclaimer:} The results in this section are preliminary. They were measured
on a single hardware configuration and should not be cited as definitive performance
numbers. We present them to indicate the general performance characteristics of the
system. Independent verification is encouraged.

\subsection{Experimental Setup}

All GPU benchmarks were conducted on an NVIDIA A100~80\,GB SXM4 with
Qwen2.5-1.5B as the model. Both rvLLM and Python vLLM~0.18.0 were tested on
identical hardware with identical prompts. CPU benchmarks (sampling, logit processing)
were measured on both Apple M5 (10 cores) and the A100 host Xeon using
\texttt{criterion} (Rust) and \texttt{timeit} (Python).

\subsection{GPU Inference}

Table~\ref{tab:gpu_bench} presents verified resource measurements. Throughput and
latency comparisons against Python vLLM are being refreshed in a follow-up
benchmark pass using a reproducible script (\texttt{bench/run.sh}).

\begin{table}[t]
  \centering
  \caption{Verified resource measurements. A100 80\,GB SXM4, Qwen2.5-1.5B.}
  \label{tab:gpu_bench}
  \small
  \begin{tabular}{@{}lr@{}}
    \toprule
    \rowcolor{HeaderBlue} \thead{Metric} & \thead{rvLLM (Rust)} \\
    \midrule
    Startup time (s)         & $\sim$7 \\
    \rowcolor{gray!8}
    Binary size              & 15\,MB \\
    CPU memory RSS           & 333\,MB \\
    \rowcolor{gray!8}
    GPU VRAM                 & 6{,}357\,MiB \\
    Output quality           & Coherent (5/5 prompts) \\
    \rowcolor{gray!8}
    Throughput (tok/s)       & \emph{Pending} \\
    P50/P95 latency          & \emph{Pending} \\
    \bottomrule
  \end{tabular}
\end{table}

\subsection{CPU Component Benchmarks}

Table~\ref{tab:cpu_bench} shows microbenchmarks for operations that run on CPU between
GPU forward passes. These are measured with \texttt{criterion} (Rust, 100~iterations,
warm cache) and \texttt{timeit} (Python/NumPy, 100~iterations).

\begin{table}[t]
  \centering
  \caption{CPU component benchmarks: sampling and logit processing. Measured on Apple M5 (10 cores, 24\,GB).}
  \label{tab:cpu_bench}
  \small
  \begin{tabular}{@{}lrrrp{3.2cm}@{}}
    \toprule
    \rowcolor{HeaderBlue} \thead{Operation} & \thead{Rust} & \thead{Python} & \thead{Speedup} & \thead{Notes} \\
    \midrule
    Combined penalties      & 2.6\,$\mu$s  & 63\,$\mu$s  & 24$\times$  & Zero-alloc iteration \\
    \rowcolor{gray!8}
    Repetition penalty (2K) & 3.1\,$\mu$s  & 34\,$\mu$s  & 11$\times$  & In-place mutation \\
    Multinomial (32K vocab)  & 12\,$\mu$s   & 66\,$\mu$s  & 5.5$\times$ & Cumulative sum + early exit \\
    \rowcolor{gray!8}
    Top-$p$ nucleus (128K)  & 1.6\,ms      & 6.9\,ms     & 4.3$\times$ & Partial sort + threshold \\
    Q4 dequant (10M elem)   & 7.1\,ms      & 9.7\,ms     & 1.4$\times$ & Chunk autovectorisation \\
    \rowcolor{gray!8}
    Batch sampling (64 seqs) & 4.3\,ms     & 36.4\,ms    & 8.5$\times$ & Rayon across 10 cores \\
    \bottomrule
  \end{tabular}
\end{table}

The most significant speedups occur in penalty computation (24$\times$) where Rust's
zero-allocation iteration over token histories avoids Python object creation overhead,
and in batch sampling (8.5$\times$) where Rayon parallelism across all cores eliminates
the GIL bottleneck.

\subsection{Limitations}

Several factors should be considered when interpreting these results:
\begin{itemize}
  \item Benchmarks use a single relatively small model (Qwen2.5-1.5B). Performance
        characteristics may differ for larger models where GPU computation dominates.
  \item Python vLLM includes optimisations (continuous batching, chunked prefill)
        that are more impactful at high concurrency. Our benchmarks do not yet cover
        high-concurrency regimes thoroughly.
  \item rvLLM does not yet implement the full set of optimisations present in mature
        Python vLLM (e.g., prefix caching is implemented but not yet benchmarked;
        CUDA graph replay is feature-gated).
  \item Throughput ratios marked ``TBD'' in Table~\ref{tab:gpu_bench} are pending
        verification by an independent benchmark agent.
\end{itemize}

% ============================================================================
\section{Related Work}
\label{sec:related}
% ============================================================================

\textbf{vLLM}~\cite{kwon2023vllm} introduced PagedAttention and demonstrated that
virtual-memory-inspired KV cache management dramatically improves serving throughput.
rvLLM reimplements the core vLLM architecture (paged attention, continuous batching,
block management) in Rust while maintaining API compatibility. The key differences are
language-level: Rust's ownership model, absence of GIL, and deterministic memory
management versus Python's ecosystem maturity and ease of modification.

\textbf{TensorRT-LLM}~\cite{tensorrtllm2024} is NVIDIA's optimised inference solution,
built on TensorRT with C++ and Python components. It achieves state-of-the-art
performance through operator fusion, quantisation, and hardware-specific kernel tuning.
However, it is tied to NVIDIA's proprietary runtime, requires model-specific compilation
steps, and has a complex build process. rvLLM targets similar performance goals with a
simpler deployment model (single binary) and full source availability.

\textbf{llama.cpp}~\cite{llamacpp2023} pioneered efficient LLM inference in C/C++,
with extensive quantisation support (GGML/GGUF formats) and broad hardware
compatibility including CPU, CUDA, Metal, Vulkan, and SYCL backends. llama.cpp excels
at single-user local inference. rvLLM targets the multi-tenant server use case with
continuous batching, concurrent request handling, and an OpenAI-compatible API layer.

\textbf{candle}~\cite{candle2023} is a Rust ML framework from Hugging Face providing
tensor operations and model implementations. While candle provides foundational tensor
primitives, rvLLM implements the full serving stack: scheduling, batching, paged
attention, streaming API, and observability. The two projects are complementary---rvLLM
could potentially use candle's tensor abstractions, though it currently uses direct
cuBLAS and custom kernel calls for maximum control.

\textbf{FlashAttention}~\cite{dao2022flashattention,dao2023flashattention2} provides
IO-aware exact attention with reduced memory footprint. rvLLM includes both a CPU
reference FlashAttention implementation for testing and a CUDA kernel for production
use.

% ============================================================================
\section{Conclusion and Future Work}
\label{sec:conclusion}
% ============================================================================

rvLLM demonstrates that a systems-language reimplementation of LLM serving
infrastructure can yield practical benefits in deployment size (15\,MB binary),
startup time ($\sim$7\,s), and memory footprint (333\,MB RSS), while maintaining
API compatibility with the existing vLLM ecosystem. The 22-crate modular architecture
enables independent development and testing of each layer, and the mock-GPU feature
makes the entire non-GPU stack accessible to developers without NVIDIA hardware.

Several directions for future work remain:

\textbf{LoRA adapter hot-swapping.} Supporting multiple LoRA adapters served from a
single base model, with runtime adapter loading and unloading, is a high-priority
feature for multi-tenant deployments.

\textbf{Vision-language models.} Extending the model runner to support multi-modal
architectures (image encoders, cross-attention) would expand the set of servable
models.

\textbf{Pipeline parallelism.} The current tensor parallelism support (NCCL bindings,
column/row parallel) handles single-node multi-GPU. Pipeline parallelism would enable
serving models that exceed single-node memory across multiple machines.

\textbf{CUDA graph integration.} Wiring the existing CUDA graph capture/replay
infrastructure into the forward pass, with careful handling of variable batch sizes,
could reduce kernel launch overhead at small batch sizes.

\textbf{Production hardening.} Fuzz testing of the API layer, load testing at 1,000+
concurrent connections, and long-running stability tests are needed before production
deployment recommendations.

The project is open-source under the Apache-2.0 license and accepts contributions for
all of the above areas.

% ============================================================================
% References
% ============================================================================
\bibliographystyle{plainnat}
\bibliography{references}

\end{document}
